

\section{Experiments}

In this section, we conduct experiments across a variety of IL problem domains: single/cross-domain IL, online/offline IL, and LfD/LfO IL settings. 
%where single-domain IL means that the learning MDP is consistent with the expert-data-collecting environment, cross-domain IL means the learning MDP and the expert-data MDP have different transition dynamics, online IL means the we can rollout the learning policy in a interative environment, offline IL means we have a xxx. LfD means the expert-data consists of state-action-next-state transitions, LfO means the expert-data contains only state-next-state transitions, 
By arranging and combining these IL domains, we obtain 8 IL tasks in all: \textit{S-on-LfD}, \textit{S-on-LfO}, \textit{S-off-LfD}, \textit{S-off-LfO}, \textit{C-on-LfD}, \textit{C-on-LfO}, \textit{C-off-LfD}, and \textit{C-off-LfO}, where S/C denotes single/cross-domain IL, on/off denotes online/offline IL, and LfD/LfO denote learning from demonstrations/observations respectively. Moreover, we also verify the scalability of CEIL on the challenging one-shot IL setting. 


Our experiments are conducted in four popular MuJoCo environments: Hopper-v2 (Hop.), HalfCheetah-v2 (Hal.), Walker2d-v2 (Wal.), and Ant-v2 (Ant.). 
In the single-domain IL setting, we train a SAC policy in each environment and use the learned expert policy to collect expert trajectories (demonstrations/observations). 
To investigate the cross-domain IL setting, we assume the two domains (learning MDP and the expert-data collecting MDP) have the same state space and action space, while they have different transition dynamics. To achieve this, we modify the torso length of the MuJoCo agents (see details in Appendix). Then, for each modified agent, we train a separate expert policy and collect expert trajectories. 
For the offline IL setting, we directly take the reward-free D4RL~\citep{fu2020d4rl} as the offline dataset, replacing the online rollout experience in the online IL setting. 



\begin{figure}[t]
	%	\vspace{-2pt}
	\small
	\begin{center}
		\includegraphics[scale=0.365]{figures/lfd/online_lfd_hopper_ant.pdf} \ 
		\includegraphics[scale=0.365]{figures/lfo/online_lfo_hopper_ant.pdf}
		% \includegraphics[scale=0.170]{figures/none.pdf}
		% \includegraphics[scale=0.170]{figures/none.pdf}
		\includegraphics[scale=0.365]{figures/lfd_cd/online_lfd_cd_hopper_ant.pdf} \ 
		\includegraphics[scale=0.365]{figures/lfo_cd/online_lfo_cd_hopper_ant.pdf} 
		\includegraphics[scale=0.365]{figures/lfd/online_lfd_hopper_ant_legend.pdf}
		\includegraphics[scale=0.365]{figures/lfo/online_lfo_hopper_ant_legend.pdf} \quad \quad \quad  
	\end{center}
	\vspace{-12pt}
	\caption{{\small Return curves on 4 online IL settings: (\textbf{a}) \textit{S-on-LfD}, (\textbf{b}) \textit{S-on-LfO}, (\textbf{c}) \textit{C-on-LfD}, and (\textbf{d}) \textit{C-on-LfO}, where the shaded area represents a 95\% confidence interval over 30 trails.
    Note that baselines cannot be applied to all the IL task settings, thus we only provide comparisons with compatible baselines (two separate legends).  % We provide a separate legend for each setting to distinguish between different baselines.
    }}
	\label{fig:online-il-results} 
	\vspace{-2pt}
\end{figure}







\begin{table}[htbp]
  \centering
  \caption{{\small Normalized scores (averaged over 30 trails for each task) on 4 offline IL settings: \textit{S-off-LfD}, \textit{S-off-LfO}, \textit{C-off-LfD}, and \textit{C-off-LfO}. Scores within two points of the maximum score are highlighted }}
  \vspace{-5pt}
  \resizebox{\linewidth}{!}{
    \begin{tabular}{clrrrrrrrrrrrrr}
    \toprule
    \multicolumn{2}{c}{\multirow{2}[2]{*}{Offline IL settings}} & \multicolumn{3}{c}{Hopper-v2} & \multicolumn{3}{c}{Halfcheetah-v2} & \multicolumn{3}{c}{Walker2d-v2} & \multicolumn{3}{c}{Ant-v2} &\multicolumn{1}{c}{\multirow{2}[2]{*}{sum}}\\
    \cmidrule(lr){3-5} \cmidrule(lr){6-8} \cmidrule(lr){9-11} \cmidrule(lr){12-14} 
    \multicolumn{2}{c}{} & \multicolumn{1}{c}{m} & \multicolumn{1}{c}{mr} & \multicolumn{1}{c}{me} & \multicolumn{1}{c}{m} & \multicolumn{1}{c}{mr} & \multicolumn{1}{c}{me} & \multicolumn{1}{c}{m} & \multicolumn{1}{c}{mr} & \multicolumn{1}{c}{me} & \multicolumn{1}{c}{m} & \multicolumn{1}{c}{mr} & \multicolumn{1}{c}{me} &  \\
    \midrule
    \multirow{7}[2]{*}{\rotatebox{90}{S-LfD}} & ORIL (TD3+BC)  & 50.9  & 22.1  & 72.7  & \cellcolor{cyan!19}44.7  & 30.2  & \cellcolor{cyan!19}87.5  & 47.1  & 26.7  & 102.6  & 46.5  & 31.4  & 61.9  & 624.3  \\
          & SQIL (TD3+BC)  & 32.6  & 60.6  & 25.5  & 13.2  & 25.3  & 14.4  & 25.6  & 15.6  & 8.0   & 63.6  & 58.4  & 44.3  & 387.1  \\
          & IQ-Learn & 21.3  & 19.9  & 24.9  & 5.0   & 7.5   & 7.5   & 22.3  & 19.6  & 18.5  & 38.4  & 24.3  & 55.3  & 264.5  \\
          & ValueDICE & 73.8  & 83.6  & 50.8  & 1.9   & 2.4   & 3.2   & 24.6  & 26.4  & 44.1  & 79.1  & 82.4  & 75.2  & 547.5  \\
          & DemoDICE & 54.8  & 32.7  & 65.4  & \cellcolor{cyan!19}42.8  & \cellcolor{cyan!19}37.0  & 55.6  & 68.1  & 39.7  & 95.0  & 85.6  & 69.0  & 108.8  & 754.6  \\
          & SMODICE & 56.1  & 28.7  & 68.0  & \cellcolor{cyan!19}42.7  & \cellcolor{cyan!19}37.7  & 66.9  & 66.2  & 40.7  & 58.2  & 87.4  & 69.9  & \cellcolor{cyan!19}113.4  & 735.9  \\
          & \textbf{CEIL (ours)} & \cellcolor{cyan!19}110.4  & \cellcolor{cyan!19}103.0  & \cellcolor{cyan!19}106.8  & 40.0  & 30.3  & 63.9  & \cellcolor{cyan!19}118.6  & \cellcolor{cyan!19}110.8  & \cellcolor{cyan!19}117.0  & \cellcolor{cyan!19}126.3  & \cellcolor{cyan!19}122.0  & \cellcolor{cyan!19}114.3  &\cellcolor{cyan!19}\textbf{1163.5}\\
    \midrule
    \multirow{3}[2]{*}{\rotatebox{90}{S-LfO}} & ORIL (TD3+BC)  & 43.4  & 25.7  & 73.0  & \cellcolor{cyan!19}44.9  & 2.4   & \cellcolor{cyan!19}81.8  & 58.9  & 16.8  & 78.2  & 33.7  & 29.6  & 67.1  & 555.4  \\
          & SMODICE & \cellcolor{cyan!19}54.5  & 26.4  & 73.7  & 42.7  & 37.9  & 66.2  & 60.6  & \cellcolor{cyan!19}38.5  & 70.9  & 85.7  & \cellcolor{cyan!19}68.3  & 116.3  & 741.7  \\
          & \textbf{CEIL (ours)} & \cellcolor{cyan!19}54.2  & \cellcolor{cyan!19}51.4  & \cellcolor{cyan!19}90.4  & \cellcolor{cyan!19}43.5  & \cellcolor{cyan!19}40.1  & 47.7  & \cellcolor{cyan!19}78.5  & 20.5  & \cellcolor{cyan!19}110.0  & \cellcolor{cyan!19}97.0  & \cellcolor{cyan!19}67.8  & \cellcolor{cyan!19}120.5  &\cellcolor{cyan!19}\textbf{821.7}\\
    \midrule
    \multirow{7}[2]{*}{\rotatebox{90}{C-LfD}} & ORIL (TD3+BC)  & 52.8  & 27.6  & 46.5  & 38.3  & 8.0   & \cellcolor{cyan!19}74.0  & 25.3  & 28.4  & 26.3  & 26.0  & 17.6  & 11.9  & 382.6  \\
          & SQIL (TD3+BC)  & 34.4  & 19.1  & 11.4  & 19.2  & 25.1  & 19.9  & 15.8  & 16.5  & 8.8   & 21.8  & 23.2  & 21.2  & 236.2  \\
          & IQ-Lean & 37.3  & 35.4  & 25.9  & 27.4  & 27.1  & 31.2  & 27.7  & 22.2  & 31.7  & 63.7  & 63.3  & 55.8  & 448.8  \\
          & ValueDICE & 22.0  & 18.3  & 18.9  & 14.0  & 11.7  & 8.7   & 11.5  & 10.0  & 8.6   & 24.1  & 21.4  & 19.2  & 188.4  \\
          & DemoDICE & 52.9  & 15.2  & 77.2  & 42.8  & \cellcolor{cyan!19}38.9  & 53.8  & 58.4  & 26.4  & 77.8  & 87.8  & 69.3  & \cellcolor{cyan!19}114.9  & 715.6  \\
          & SMODICE & 55.4  & 21.4  & 71.2  & \cellcolor{cyan!19}42.7  & \cellcolor{cyan!19}38.0  & 64.6  & 68.4  & \cellcolor{cyan!19}34.2  & \cellcolor{cyan!19}80.4  & 87.4  & 70.4  & \cellcolor{cyan!19}115.7  & 749.7  \\
          & \textbf{CEIL (ours)} & \cellcolor{cyan!19}58.4  & \cellcolor{cyan!19}39.8  & \cellcolor{cyan!19}81.6  & \cellcolor{cyan!19}42.6  & \cellcolor{cyan!19}38.3  & 46.6  & \cellcolor{cyan!19}76.5  & 21.1  & \cellcolor{cyan!19}81.1  & \cellcolor{cyan!19}91.6  & \cellcolor{cyan!19}88.0  & \cellcolor{cyan!19}115.3  &\cellcolor{cyan!19}\textbf{780.9}\\
    \midrule
    \multirow{3}[2]{*}{\rotatebox{90}{C-LfO}} & ORIL (TD3+BC)  & \cellcolor{cyan!19}55.5  & 18.2  & 55.5  & 40.6  & 2.9 & \cellcolor{cyan!19}73.0  & 26.9  & 19.4  & 22.7  & 11.2  & 21.3  & 10.8  & 358.0  \\
          & SMODICE & \cellcolor{cyan!19}53.7  & 18.3  & \cellcolor{cyan!19}64.2  & \cellcolor{cyan!19}42.6  & \cellcolor{cyan!19}38.0  & 63.0  & 68.9  & \cellcolor{cyan!19}37.5  & 60.7  & 87.5  & \cellcolor{cyan!19}75.1  & \cellcolor{cyan!19}115.0  & \cellcolor{cyan!19}\textbf{724.4}  \\
          & \textbf{CEIL (ours)} & 44.7  & \cellcolor{cyan!19}44.2  & 48.2  & \cellcolor{cyan!19}42.4  & \cellcolor{cyan!19}36.5  & 46.9  & \cellcolor{cyan!19}76.2  & 31.7  & \cellcolor{cyan!19}77.0  & \cellcolor{cyan!19}95.9  & 71.0  & 112.7  &\cellcolor{cyan!19}\textbf{727.3}\\
    \bottomrule
    \end{tabular}%
    }
  \label{tab:offline-il-results}%
  \vspace{-5pt}
\end{table}%











% SAIL~\citep{liu2019state}, 



\subsection{Evaluation Results}

To demonstrate the versatility of the CEIL idea, we collect 20 expert trajectories (demonstrations in LfD or observations in LfO) for each environment and compare CEIL to GAIL~\citep{ho2016generative}, AIRL~\citep{fu2017learning}, SQIL~\citep{reddy2019sqil}, IQ-Learn~\citep{garg2021iq}, ValueDICE~\citep{kostrikov2019imitation}, GAIfO~\citep{torabi2018generative},  ORIL~\citep{zolna2020offline}, DemoDICE~\citep{kim2022demodice},  SMODICE~\citep{ma2022smodice} (see details in Appendix).
Note that these baseline methods cannot be applied to all the IL task settings, thus we only provide comparisons with compatible baselines in each IL setting. 





\textbf{Online IL.} 
In Figure~\ref{fig:online-il-results}, we provide the return (cumulative rewards) curves of our method and baselines on 4 online IL settings: \textit{S-on-LfD} (\textit{top-left}), \textit{S-on-LfO} (\textit{top-right}), \textit{C-on-LfD} (\textit{bottom-left}), and \textit{C-on-LfO} (\textit{bottom-right}) settings. 
As can be seen, CEIL quickly achieves expert-level performance in \textit{S-on-LfD}. 
When extended to \textit{S-on-LfO}, CEIL also yields better sample efficiency compared to baselines. 
Further, considering the complex cross-domain setting, we can see those baselines SQIL and IQ-Learn (in \textit{C-on-LfD} and \textit{C-on-LfO}) suffer from the domain mismatch, leading to performance degradation at late stages of training, while CEIL can still achieve robust performance.







\begin{figure}[H]
	\vspace{-5pt}
	\begin{center}
		\includegraphics[scale=0.335]{figures/num/all_Hopper-v2_.pdf}
		\includegraphics[scale=0.335]{figures/num/all_Ant-v2_.pdf} \ \ 
		\includegraphics[scale=0.345]{figures/window/Hopper-v2_online_lfd.pdf}
		\includegraphics[scale=0.345]{figures/window/Ant-v2_online_lfd.pdf}
		\includegraphics[width=0.95\columnwidth]{figures/num/all_both_legend.pdf}
	\end{center}
	\vspace{-11pt}
	\caption{{\small Ablating (\textbf{a}, \textbf{b}) the number of  expert demonstrations and (\textbf{c}, \textbf{d}) the trajectory window size.}}
	\label{fig:varying_num_window} 
	\vspace{-8pt}
\end{figure}




\textbf{Offline IL.} 
Next, we evaluate CEIL on the other 4 offline IL settings: \textit{S-off-LfD}, \textit{S-off-LfO}, \textit{C-off-LfD}, and \textit{C-off-LfO}. In Table~\ref{tab:offline-il-results}, we provide the normalized return of our method and baseline~methods~on 
reward-free D4RL~\citep{fu2020d4rl}
 medium (m), medium-replay (mr), and medium-expert (me) datasets. 
We can observe that CEIL achieves a significant improvement over the baseline methods in both \textit{S-off-LfD} and \textit{S-off-LfO} settings. Compared to the state-of-the-art offline baselines, CEIL also shows competitive results on the challenging cross-domain offline IL settings (\textit{C-off-LfD} and \textit{C-off-LfO}). 
% In Appendix I, we provide more experimental results across single/cross-domain IL settings, which also demonstrate the competitiveness of CEIL compared to state-of-the-art offline IL baselines. 





\textbf{One-shot IL.} Then, we explore CEIL on the one-shot IL tasks, where we expect CEIL can adapt its behavior to new IL tasks given only one trajectory for each task (mismatched MDP,~see~Appendix). 


\begin{wraptable}{r}{0.42\textwidth}
	\centering
    \vspace{-10pt}
	\resizebox{\linewidth}{!}{
    \begin{tabular}{llrrrr}
    \toprule
    \multicolumn{2}{l}{\textbf{One-shot IL}} & \multicolumn{1}{c}{Hop.} & \multicolumn{1}{c}{Hal.} & \multicolumn{1}{c}{Wal.} & \multicolumn{1}{c}{Ant.} \\ % & \multicolumn{1}{c}{sum}\\
    \midrule
    \multirow{3}[2]{*}{\rotatebox{90}{Online}} & SQIL   & 16.8  & 1.1   & 3.5   & 4.2 \\ %& 25.5  \\
          & IQ-Learn & 4.6   & 0.2   & 1.7   & 7.5 \\ %& 14.0  \\
          & \textbf{CEIL} (LfD)  & \cellcolor{cyan!19}29.9  &  \cellcolor{cyan!19}2.5  & \cellcolor{cyan!19}31.7  & 20.5 \\ %& \cellcolor{cyan!19}\textbf{84.5} \\
          & \textbf{CEIL} (LfO)  & 17.8  &  \cellcolor{cyan!19}3.2  &  5.6  & \cellcolor{cyan!19}29.7 \\ %& 56.3 \\
    \midrule
    \multirow{6}[2]{*}{\rotatebox{90}{Offine}} & ORIL  & 14.7  & 0.2   & 6.9   & 17.4 \\ %& 39.1  \\
          & SQIL  & 7.4   & 0.8   & 4.6   & 12.5 \\ %& 25.3  \\
          & IQ-Learn & 18.8  & 1.2   & 4.0   & 19.3 \\ %& 43.3  \\
          & DemoDICE   & 76.5  & -0.5  & -0.1  & 19.5 \\ %& 95.4  \\
          & SMODICE & 78.0  & 1.1   & 8.1   & \cellcolor{cyan!19}24.6 \\ %& 111.9\\
          & \textbf{CEIL} (LfD)  & \cellcolor{cyan!19}85.6  &  \cellcolor{cyan!19}5.6   & \cellcolor{cyan!19}67.1  & \cellcolor{cyan!19}24.3 \\ %& \cellcolor{cyan!19}\textbf{182.6} \\
          & \textbf{CEIL} (LfO)  & 72.2  &  \cellcolor{cyan!19}5.1   & \cellcolor{cyan!19}70.0  &  19.4 \\ %& 171.8 \\
    \bottomrule
    \end{tabular}%
    }
    \vspace{-5pt}
    \caption{{\small Normalized results on one-shot IL, where CEIL shows prominent transferability.}}
  \label{tab:one-shot-il-results}
  \vspace{-10pt}
\end{wraptable}
We first pre-train an embedding function and a contextual policy in the training domain (online/offline IL), then infer a new contextual variable and evaluate it on the new task. To facilitate comparison to baselines, we similarly pre-train a policy network (using baselines) and run BC on top of the pre-trained policy by using the provided demonstration. 
Consequently, such a baseline+BC procedure cannot be applied to the (one-shot) LfO tasks. 
% In , we provide the results in 4 Mujoco tasks. 
The results in Table~\ref{tab:one-shot-il-results} show that baseline+BC struggles to transfer their expertise to new tasks. 
Benefiting from the hindsight framework, CEIL shows better one-shot transfer learning performance on 7 out of 8 one-shot LfD tasks and retains higher scalability and generality for both one-shot LfD and LfO IL tasks. 
% % On the other hand, the DICE-based methods, DemoDICE and SMODICE, achieve higher returns, especially on Hopper and Ant, which is consistent with the advantage of versatility discussed and shown in ~\citep{ma2022smodice}.


\vspace{-1pt}
\subsection{Analysis of CEIL}
\vspace{-2pt}

\begin{wraptable}{r}{0.58\textwidth}
    \vspace{-10pt}
	\centering
	\resizebox{\linewidth}{!}{
	\begin{tabular}{lrrrr}
		\toprule
            \textbf{Hybrid offline IL settings} & \multicolumn{1}{l}{Hop.} & \multicolumn{1}{l}{Hal.} & \multicolumn{1}{l}{Wal.} & \multicolumn{1}{l}{Ant.} \\
		\midrule
		S-LfD & \cellcolor{cyan!8}29.4  & \cellcolor{cyan!8}69.9  & \cellcolor{cyan!8}42.8  & \cellcolor{cyan!8}84.9  \\
		S-LfD + S-LfO & \cellcolor{cyan!20}30.4  & \cellcolor{cyan!10}68.6  & \cellcolor{cyan!20}42.3  & \cellcolor{cyan!30}91.6  \\
		S-LfD + S-LfO + C-LfD & \cellcolor{cyan!30}30.7  & \cellcolor{cyan!30}71.7  & \cellcolor{cyan!30}42.9  & \cellcolor{cyan!20}89.2  \\
		S-LfD + S-LfO + C-LfD + C-LfO & \cellcolor{cyan!50}58.6  & \cellcolor{cyan!50}79.6  & \cellcolor{cyan!50}43.7  & \cellcolor{cyan!50}98.0  \\
		\bottomrule
	\end{tabular}
	}
    \vspace{-5pt}
	\caption{{\small The normalized results of CEIL, showing that CEIL can consistently digest useful (task-relevant) information and boost its performance, even under a hybrid of offline IL settings.}}
	\label{tab:ceil-hybrid}
    \vspace{-5pt}
\end{wraptable}
\textbf{Hybrid IL settings.} In real-world, many IL tasks do not correspond to one specific IL setting, and instead consist of a hybrid of several IL settings, each of which passes a portion of task-relevant information to the IL agent. For example, we can provide the agent with both demonstrations and state-only observations and, in some cases, cross-domain demonstrations (S-LfD+S-LfO+C-LfD). To examine the versatility of CEIL, we collect a separate expert trajectory for each of the four offline IL settings, and study CEIL's performance under hybrid IL settings. As shown in Table~\ref{tab:ceil-hybrid}, we can see that by adding new expert behaviors on top of LfD, even when carrying relatively less~supervision (\eg actions are absent in LfO), CEIL can still improve the performance. % by increasing hybrid behaviors.




\textbf{Varying the number of demonstrations.} In Figure~\ref{fig:varying_num_window} (a, b), we study the effect of the number of expert demonstrations on CEIL's performance. Empirically, we reduce the number of training demonstrations from 20 to 1, and report the normalized returns at $1 \text{M}$ training steps. We can observe that across both online and offline (D4RL *-medium) IL settings, CEIL shows more robust performance with respect to different numbers of demonstrations compared to baseline methods. 

\textbf{Varying the window size of trajectory.} Next we assess the effect of the trajectory window size (\ie the length of trajectory $\btau$ used for the embedding function $f_\phi$ in Equation~\ref{eq:traj-self-consistent}). In Figure~\ref{fig:varying_num_window} (b, c), we ablate the number of the window size in 4 LfD IL instantiations. We can see that across a range of window sizes, CEIL remains stable and achieves expert-level performance. 

 



